\section{Conclusion}

The United States interstate highway system covers 79.6\% of the national population within the 30-mile access standard. That leaves 66.5 million Americans in 1,510 counties beyond the threshold --- a gap that is not random but concentrated in four identifiable geographic zones: the Northern Tier, the Appalachians, the Rural Gulf South, and the Rural West.

These are not marginal populations. They include the agricultural heartland of Montana and North Dakota, the coalfield communities of Appalachia, the bayou parishes of Louisiana, the rural ranching counties of the Texas interior. They are geographically isolated by the absence of investment decisions made 70 years ago, when the interstate system was designed to connect existing cities rather than to achieve minimum access for the full territory.

Three findings are the paper's primary contributions.

\textbf{The gap is 66.5 million people.} Population-weighted coverage analysis from Census county centroids establishes, for the first time with public data, the exact population without adequate interstate access. This is a plannable, closable gap. It is not a statement about the impossibility of coverage; it is a target.

\textbf{Twelve corridors close most of it.} The top 12 proposed corridors, ranked by coverage efficiency (gap population served per \$B of investment), would bring 11.1 million gap-county residents within threshold --- raising national coverage from 79.6\% to 83\%. The top four (I-3, Northern Tier, I-14, I-69) serve the four gap zones with the highest efficiency ratios. Rural access spurs address the most isolated Type 3 communities.

\textbf{The coverage gap and the economic opportunity gap are the same gap.} Gap counties score 41\% higher on the C3 (Economic Opportunity Access) dimension than non-gap counties, with mean buffer GDP per capita at 77\% of the national average. The investment that closes the coverage gap also connects below-average economic regions to the national economy. The transportation case and the economic case point to the same corridors.

\subsection*{Policy Implications}

The priority corridor ranking has direct implications for Interstate 2.0 investment sequencing. I-3 and the Northern Tier are not politically visible corridors --- they do not serve large metropolitan constituencies and they do not appear in ATRI's top congestion bottlenecks. Under a volume-based investment criterion, they would be indefinitely deferred. Under a coverage-and-economic-opportunity criterion, they are the highest-priority new construction in the national network.

This is the policy choice Interstate 2.0 makes explicit: the managed freight lane investment in Atlanta and Chicago addresses where the system is most stressed; the missing link investment in Appalachia and the Northern Plains addresses where the system never arrived. Both are necessary. The tier-based investment framework of \citet{ROUTE_E2} provides the sequencing mechanism.

\subsection*{Future Work}

This paper's coverage analysis uses county centroids, which overstate access deficits in large geographic counties and understate them in counties with strong population clustering away from the centroid. A tract-level coverage analysis using Census block group centroids would provide substantially more precision, particularly for California, Arizona, and Nevada. Such an analysis is computationally feasible with the ROUTE system but requires tract-level population data not yet joined in the current pipeline.

The coverage analysis measures access to the \textit{existing} interstate network; it does not measure access to the proposed T2 and T3 improvements that I2.0 includes. Paper B.3 \citep{ROUTE_B3} extends the coverage analysis to climate exposure (which counties are within the 30-mile threshold but on corridors with high flood vulnerability), and paper E.2 \citep{ROUTE_E2} integrates coverage, congestion, resilience, and transit access into a unified investment priority score.
